\documentclass[11pt]{article}

% ---------------------------------------------------------------------------
% The F2Z verifier note. Build: latexmk -pdf main.tex
% Describes verify_mle_eval_mod_q_ligerito as implemented in this repository,
% treating the flock Ligerito verifier as a black-box oracle.
% ---------------------------------------------------------------------------

\usepackage[margin=1.05in]{geometry}
\usepackage{amsmath,amssymb,amsthm}
\usepackage{booktabs}
\usepackage{array}
\usepackage{enumitem}
\usepackage{microtype}
\usepackage{xcolor}
\usepackage[colorlinks=true,linkcolor=blue!50!black,citecolor=blue!50!black,urlcolor=blue!50!black]{hyperref}

% --- notation (aligned with the upstream f2-int-eval note) ------------------
\newcommand{\F}{\mathbb{F}}
\newcommand{\Z}{\mathbb{Z}}
\newcommand{\K}{\mathbb{K}}            % the binding field GF(2^128)
\newcommand{\Kx}{\K^{\times}}
\newcommand{\Ftwo}{\F_2}
\newcommand{\Fq}{\F_q}                 % the evaluation field (char != 2)
\newcommand{\bits}{\{0,1\}}
\newcommand{\INT}{\mathsf{INT}}
\newcommand{\mle}[1]{\widetilde{#1}}
\newcommand{\eq}{\mathsf{eq}}
\newcommand{\cw}{c_w}
\newcommand{\albert}[1]{\textcolor{teal}{Albert: {#1}}}
% \code: typewriter names; underscores allow a line break after them, so
% long function names wrap instead of overflowing the margin.
\newcommand{\code}[1]{{\def\_{\textunderscore\allowbreak}\texttt{#1}}}
\newcommand{\RecVerify}{\mathsf{LigVerify}}

% Typewriter-heavy paragraphs: cmtt spaces have no stretch, so give the
% line breaker slack instead of letting it emit overfull lines.
\tolerance=1600
\setlength{\emergencystretch}{1.8em}
\hbadness=2500

\newtheorem{definition}{Definition}
\newtheorem{remark}{Remark}
\newtheorem{obligation}{Obligation}

\newcommand{\Vc}[1]{\textbf{[V#1]}}    % verifier check anchors

\title{The F2Z verifier for the mod-$q$ opening,\\
  with Ligerito as a black box}
\author{The \code{f2z-pcs} repository}
\date{2026-07-16}

\begin{document}
\maketitle

\begin{abstract}
\noindent
This note specifies, at code-review granularity, the verifier of the F2Z
polynomial commitment scheme for its core statement — the mod-$q$
integer-MLE evaluation
$\smash{\mle{\INT(D)}}(r) = y \in \Fq$ over an $\Ftwo$-committed bit
matrix — as implemented by \code{verify\_mle\_eval\_mod\_q\_ligerito} in
\code{src/ligerito\_flock.rs}.  The flock Ligerito proximity verifier is
treated throughout as a correct black box with a one-line interface
contract (Definition~\ref{def:oracle}); its internals are explicitly out
of scope.  The note walks the soundness chain in verification order —
committed bits $\to$ (opening) residual bit-MLE claims $\to$ (forest GKR)
roots $\alpha^{u}$ $\to$ (generator binding) integers $u$ $\to$ (chunk
recombination) $y \in \Fq$ — and for every verifier-side quantity states
whether it is \emph{recomputed from public data}, \emph{bound by a
transcript absorb}, or \emph{oracle-certified}.  A reader holding this
note and the source should be able to confirm, function by function, that
every stated check exists and that no stated check is silently missing.
\end{abstract}

\tableofcontents

% ===========================================================================
\section{Scope, review contract, and the oracle}
\label{sec:scope}

F2Z\footnote{Repository \code{f2z-pcs}; a standalone extraction of the
most-optimized \code{f2-int} pipeline from the \code{zinc-plus} repository
(branch \code{f2-int-unified-ligerito}, snapshot 2026-07).  The
construction is due to Lev Soukhanov (the char2-fieldswitch note, \S9);
the as-deployed long-form exposition lives in
\code{zinc-plus}'s \code{documentation/f2-int-eval-doc/}.  This note
describes the extracted code, which has drifted from the deployment; where
the two disagree, this repository's code is authoritative for this note.}
proves an \emph{integer} multilinear-extension evaluation over data whose
\emph{bits} are committed with a characteristic-2 commitment.  The
integer row-fold is carried in the exponent of $\K = \mathrm{GF}(2^{128})$
and certified by a GKR grand-product forest; the only opener is the
flock-backed ring-switch + recursive Ligerito pipeline.

\paragraph{What this note covers.}  Exactly one entry point and its
callees:
\begin{center}
\footnotesize
\begin{tabular}{p{0.56\textwidth}p{0.36\textwidth}}
\toprule
Function(s) & File \\
\midrule
\code{verify\_mle\_eval\_mod\_q\_ligerito} & \code{src/ligerito\_flock.rs} \\
\code{verify\_int\_eval\_merged\_common} & \code{src/ligerito.rs} \\
\code{verify\_merged\_forest}, \code{mle\_at}, \code{absorb\_gfs} & \code{src/merged\_forest.rs} \\
\code{MLSumcheck::verify\_as\_subprotocol} & \code{src/piop/sumcheck.rs} \\
\code{MultiDegreeSumcheck::verify\_as\_subprotocol} & \code{src/piop/sumcheck/multi\_degree.rs} \\
\code{is\_generator}, \code{mod\_q\_chunk\_width}, \code{mod\_q\_num\_chunks},
  \code{chunk\_row\_weights}, \code{row\_bit\_weights},
  \code{recombine\_read\_off}, \code{FixedBasePow} & \code{src/pcs.rs} \\
ring-switch pieces: \code{absorb\_sv}, \code{transpose\_bits\_128},
  \code{residual\_b\_evals} & \code{src/ligerito.rs} \\
\code{IntEvalRsLigModQProof::to\_bytes} / \code{from\_bytes} &
  \code{src/ligerito\_flock.rs}, \code{src/proof\_codec.rs} \\
\code{Blake3Transcript}, \code{absorb\_slice}, \code{get\_field\_challenge} & \code{src/transcript/} \\
\code{ZincChallenger} (the flock challenger bridge) & \code{src/ligerito\_flock.rs} \\
\bottomrule
\end{tabular}
\end{center}
Line numbers drift; navigate by function name.

\paragraph{What this note does not cover.}
\begin{itemize}[itemsep=1pt]
\item \emph{The Ligerito / ring-switch proximity verifier internals}
  (\code{flock-core}).  It appears only as the oracle of
  Definition~\ref{def:oracle}.  Whenever the text says ``the opening is
  checked,'' that is the oracle call, full stop.
\item The batched multi-polynomial variants
  (\code{verify\_rs\_ligerito\_batch}) and the virtual-XOR variants
  (\code{verify\_mle\_eval\_mod\_q\_ligerito\_with\_virtual\_xors},
  \code{\ldots\_claims\_only}).  They exist as public API and are
  discussed only in the one paragraph of
  Section~\ref{sec:trust} that states their extra trust obligations.
\item The prover, the commit path, and all performance machinery (lazy
  forest schedules, NEON kernels).  These are semantics-preserving and
  pinned byte-identical by tests (Section~\ref{sec:tests}); the verifier
  never depends on which schedule produced a proof.
\end{itemize}

\paragraph{The oracle.}  The single point of contact with
\code{flock-core} on this path is the following function, used once per
verification:

\begin{definition}[Ligerito opening oracle]\label{def:oracle}
$\RecVerify(\mathit{vc},\, \pi_{\mathrm{lig}},\, m_p,\, \beta,\,
\mathsf{root},\, \mathsf{EvalB},\, \mathcal{C}) \in \{\mathrm{accept},
\mathrm{reject}\}$
— in code,
\code{flock\_core::pcs::ligerito::recursive\_verifier\_with\_basis\_succinct}.
Inputs: a verifier config $\mathit{vc}$ (derived from the instance shape,
Section~\ref{sec:config}); the flock proof object $\pi_{\mathrm{lig}}$; the
packed variable count $m_p$; a target value $\beta \in \K$; the Merkle
commitment root; a callback $\mathsf{EvalB}$ which, given a prefix
$(\rho_1,\dots,\rho_k) \in \K^k$ and a tail width $k'$ with $k + k' = m_p$,
returns the $2^{k'}$ values $\mle{B}(\rho_1,\dots,\rho_k,\mathsf{bits}(y))$
of a weight function $B \colon \bits^{m_p} \to \K$ known to the verifier;
and a Fiat–Shamir challenger $\mathcal{C}$ that continues the caller's
transcript.

To be precise about what $\mathsf{root}$ commits to: at commit time,
flock's \code{commit} RS-encodes the \emph{message} vector
$P \in \K^{2^{m_p}}$ (an interleaved additive-NTT encoding) and builds
the Merkle tree over the resulting \emph{codeword} — so the root binds
the codeword directly, and binds $P$ only through the code.  Both the
encoding and the tree live inside the box; the packing of
Section~\ref{sec:statement} specifies which $P$ the honest committer
encodes.

\emph{Assumed contract:} if the oracle accepts, then — except with
negligible probability over the challenger's randomness — $\mathsf{root}$
binds a unique message $P \in \K^{2^{m_p}}$ (uniqueness through the
Merkle tree, the code's distance, and the proximity test is exactly the
boxed guarantee), and that message satisfies
\[
  \textstyle\sum_{y \in \bits^{m_p}} P[y]\cdot B(y) \;=\; \beta ;
\]
for an honestly generated commitment, that $P$ is the committed packed
vector.  Completeness holds for honestly generated proofs.  Everything
else — the proximity test, Merkle verification, query counts, grinding,
out-of-domain sampling, and the oracle's own transcript sub-schedule —
is inside the box.
\end{definition}

Note the contract is an \emph{inner-product} (weighted-sum) opening.
The familiar point-evaluation form is the special case where $B$ is an
$\eq$-derived tensor; the mod-$q$ verifier hands the oracle an
$\eta$-batched combination of $L$ such functionals
(Section~\ref{sec:linkage}), so the inner-product form is the honest
interface.

The contract is also deliberately stated over the packed vector $P$
rather than over the bit matrix $M$ (or the data $D$): flock's verifier
is bit-agnostic — it certifies a $\K$-linear functional of the
$\K$-elements that were committed, and knows nothing about their bit
semantics.  The bits enter one level up, on the zinc side: the packing of
Section~\ref{sec:statement} makes $P \leftrightarrow M$ a bijective
reindexing (so a commitment to $P$ \emph{is} a commitment to $M$), and
the ring switch (Section~\ref{sec:linkage}) is precisely the step that
converts each bit-level claim $\mle{M}(\mathit{pt}_l) = \mu_l$ into an
inner-product claim $\sum_y P[y]\, B_l(y) = \beta_l$ of the oracle's
form.  Stating the contract over $M$ or $D$ would silently fold that
zinc-side reduction — which this note must keep reviewable — into the
black box.

\paragraph{The soundness chain.}  \code{docs/DESIGN.md} states it; the
verifier walks it in this order, one section per link:
\begin{gather*}
\underbrace{\text{committed bits}}_{\text{oracle, \S\ref{sec:linkage}}}
\;\to\;
\underbrace{\text{residual bit-MLE claims } \mle{M}(\mathit{pt}_l) = \mu_l}_{\text{ring switch, \S\ref{sec:linkage}}}
\;\to\;
\underbrace{\text{exit claims } e_d^{(l)}}_{\text{pre-sumcheck, \S\ref{sec:linkage}}}
\;\to\; \\[2pt]
\underbrace{\text{roots } \alpha^{u}}_{\text{forest, \S\ref{sec:forest}}}
\;\to\;
\underbrace{\text{integers } u}_{\text{binding, \S\ref{sec:binding}}}
\;\to\;
\underbrace{y \in \Fq}_{\text{read-off, \S\ref{sec:binding}}}
\end{gather*}
The verifier \emph{executes} the chain left-to-right per chunk in a
different operational order (range/binding checks come first,
Section~\ref{sec:schedule}); the chain above is the direction in which
trust flows once all checks pass.

\paragraph{The computing-vs-binding discipline.}  A verifier that merely
\emph{recomputes} a value from prover-chosen inputs binds nothing.
Throughout, every verifier-side quantity is classified as one of:
\begin{itemize}[itemsep=1pt]
\item[\textbf{(R)}] \emph{recomputed} from statement data (public inputs
  the caller vouches for) and verifier randomness;
\item[\textbf{(A)}] prover-supplied and \emph{bound by a transcript
  absorb} before some later challenge, and/or \emph{checked equal} to an
  already-bound quantity;
\item[\textbf{(O)}] \emph{oracle-certified} — its correctness is exactly
  Definition~\ref{def:oracle}'s guarantee.
\end{itemize}
Table~\ref{tab:provenance} classifies every quantity on the path.

% ===========================================================================
\section{Statement, instance shape, and the proof object}
\label{sec:statement}

\subsection{Instance shape and the committed matrix}

\code{IntEvalParams \{ t, s, word\_bits: W \}} (\code{src/pcs.rs}) fixes
the shape: the data $D$ is a $2^t \times 2^s$ matrix of $W$-bit cells,
$W$ a power of two, flat index $\code{cell\_index}(b,c) = (b \ll s)\,|\,c$.
Derived quantities used throughout:
\[
  d = t + \log_2 W, \qquad
  m_p = d - 7 + s, \qquad
  m = m_p + 7 = d + s
\]
(code: \code{row\_bit\_vars} and \code{packed\_vars}).
The \emph{committed object} is the single-bit matrix
$M \in \Ftwo^{2^s \times 2^d}$ with
\[
  M[c][i] = \text{bit } j \text{ of } D(b,c), \qquad
  i = (b \ll \log_2 W)\,|\,j ,
\]
i.e.\ row $c$ concatenates the bits of column $c$'s cells, cell-major,
bit-minor (\code{repack\_leaf\_bits}).  Its multilinear extension
$\mle{M} \colon \K^{d+s} \to \K$ places the $d$ row-bit coordinates
\emph{low} and the $s$ column coordinates \emph{high}, LSB-first within
each block; all $\eq$ tables in the crate follow the same LSB-first
convention (\code{build\_eq\_x\_r\_vec}).

For the commitment, the bits are packed $128$ per $\K$-element along the
low $7$ row-bit coordinates ($\code{LOG\_PACKING} = 7$): the packed vector
$P \in \K^{2^{m_p}}$ has
\[
  P[\,c \cdot 2^{\,d-7} + i_{\mathrm{hi}}\,]
  \;=\; \textstyle\sum_{v=0}^{127}\, M[c][\,(i_{\mathrm{hi}} \ll 7)\,|\,v\,]\cdot \beta_v ,
\]
where $\beta_v$ is the $v$-th bit of the $\K$-element representation
(monomial basis, LSB-first; \code{commit\_rs\_flock\_from\_rows}).  $P$ is
what flock commits: \code{flock\_core::pcs::commit::commit(\&p\_msg, \&PcsParams
\{ m: $m_p + 7$, log\_inv\_rate, log\_batch\_size, profile \})}.  Note the
Merkle root is built over the RS-\emph{encoded} (interleaved) codeword of
$P$, not over $P$ itself — $P$ is bound only through the code, which is
part of the oracle's boxed guarantee (Definition~\ref{def:oracle}); the
encoding and the tree are inside the box, and the verifier only ever
touches the \code{Commitment}'s root and shape parameters.  \emph{Precondition:} the commitment handed to the verifier
was produced at this same shape; the verifier does not cross-check
\code{commitment.params} against \code{p} (a mismatch is left to the
oracle to reject).

\subsection{The claim and the chunk decomposition}

The verifier's signature is
\begin{center}\small
\code{verify\_mle\_eval\_mod\_q\_ligerito<R>(transcript, commitment, proof, p,
row\_weights\_q, col\_weights, alpha, claimed, q\_bits, vc)}
\end{center}
with $R$ any ring of characteristic $\ne 2$ exposing
\code{From<u128>}, $+$, $\times$.  The relation verified is, for the
caller-supplied weight vectors
$w \in [0, 2^{\,q_{\mathrm{bits}}})^{2^t}$ (\code{row\_weights\_q}, as
\code{u128}s) and $w' \in R^{2^s}$ (\code{col\_weights}):
\begin{equation}\label{eq:claim}
  y \;=\; \sum_{c \in \bits^s} w'_c \sum_{b \in \bits^t} w_b \cdot \INT(D(b,c))
  \quad\text{in } R,
\end{equation}
where $\INT(D(b,c)) = \sum_j 2^j M[c][(b \ll \log_2 W)|j]$ is the cell's
integer value.  The \emph{intended instantiation} is the genuine MLE
opening: $R = \Fq$, $w_b = \eq(b, r_1) \bmod q \in [0,q)$ the canonical
integer representatives, $w'_c = \eq(c, r_2) \in \Fq$, whence
$y = \mle{\INT(D)}(r_1, r_2)$.  The verifier itself is agnostic: any
weights satisfying the width precondition are accepted, and soundness is
with respect to~\eqref{eq:claim} for the weights given.

Full-width row weights ($\approx q_{\mathrm{bits}} \approx 100$ bits)
would overflow the exponent's injectivity budget, so they are chunked
(\code{src/pcs.rs}):
\begin{gather*}
  \cw \;=\; 127 - t - W \quad (\code{mod\_q\_chunk\_width};\ \text{asserts } t + W \le 126),
  \\[1pt]
  L \;=\; \max\!\left(\lceil q_{\mathrm{bits}} / \cw \rceil,\, 1\right)
  \quad (\code{mod\_q\_num\_chunks}),
\end{gather*}
\[
  W^{(l)}_b \;=\; (w_b \gg \cw l) \wedge (2^{\cw} - 1)
  \quad (\code{chunk\_row\_weights}), \qquad l = 0, \dots, L-1,
\]
so $w_b = \sum_l W^{(l)}_b\, 2^{\cw l}$ \emph{exactly} whenever
$w_b < 2^{\cw L}$ — guaranteed by the precondition
$w_b < 2^{q_{\mathrm{bits}}} \le 2^{\cw L}$.  (Weights at or above
$2^{\cw L}$ would be silently truncated identically on both sides; the
verified relation is always \eqref{eq:claim} with $w_b$ replaced by
$\sum_l W^{(l)}_b 2^{\cw l}$.)  Each chunk's honest fold
\begin{equation}\label{eq:chunkfold}
  u^{(l)}_c \;=\; \sum_b W^{(l)}_b \cdot \INT(D(b,c))
  \;\le\; (2^{\cw}-1)(2^W-1)\,2^{t} \;<\; 2^{\,\cw + t + W} \;=\; 2^{127}
\end{equation}
fits below the generator order and binds injectively in the exponent
(Section~\ref{sec:binding}).  Prover and verifier derive
$(\cw, L, W^{(l)})$ from the same public data; nothing about the chunking
rides the proof.

\subsection{The proof object}

\code{IntEvalRsLigModQProof} (\code{src/ligerito\_flock.rs}):

\begin{center}\small
\begin{tabular}{p{0.27\textwidth}p{0.62\textwidth}}
\toprule
Field & Role \\
\midrule
\code{mfs: Vec<MergedForestProof>} &
  One merged-forest proof per chunk ($L$ total).  Each is a
  \code{Vec<MergedLayer>} of $d$ layers; a \code{MergedLayer} holds an
  optional phase-A sumcheck proof \code{sc\_x} (absent at layer $0$), the
  phase-B sumcheck proof \code{sc\_c}, and the closing child pair
  \code{pair} $\in \K^2$.
  \emph{The per-tree roots are not proof data}:
  the verifier derives them from \code{us} (Section~\ref{sec:forest}). \\
\code{us: Vec<Vec<u128>>} &
  \code{us[l]} $= (u^{(l)}_c)_{c \in \bits^s}$, the $2^s$ sent chunk folds
  of chunk $l$; claimed to equal \eqref{eq:chunkfold}. \\
\code{presums: Vec<Multi\allowbreak DegreeSumcheckProof>} &
  One de-black-boxing pre-sumcheck per chunk: $d$ rounds, one group,
  degree $2$ when honest (the degree field rides the proof;
  Remark~\ref{rem:degree}), plus the group's claimed sum. \\
\code{rings: Vec<RingSwitchProof>} &
  One ring-switch message per chunk: \code{s\_v}, the $128$ partial
  evaluations $s^{(l)}_v = \mle{M}_v(\mathit{pt}_l^{\mathrm{hi}})$
  (Section~\ref{sec:linkage}). \\
\code{lig: LigeritoProof} &
  flock's serde proof object for the single batched opening — opaque here
  (Definition~\ref{def:oracle}). \\
\bottomrule
\end{tabular}
\end{center}

\subsection{Serialization: the host proof stream}
\label{sec:serialization}

\code{to\_bytes} / \code{from\_bytes} (\code{src/ligerito\_flock.rs},
codec primitives in \code{src/proof\_codec.rs}).  All scalars are
little-endian.  The stream grammar is, in order:

\begin{center}\small
\begin{tabular}{ll}
\toprule
Bytes & Content \\
\midrule
\code{u64} & $L$ (chunk count) \\
\emph{per chunk} $l$: & \\
\quad \code{u64} & layer count ($= d$ when honest) \\
\quad \emph{per layer}: \code{u64} & \code{sc\_x} presence flag ($0$ or $1$) \\
\quad\quad [\code{Transcribable}] & \code{sc\_x} (iff flag $= 1$) \\
\quad\quad \code{Transcribable} & \code{sc\_c} \\
\quad\quad $2 \times 16$\,B & closing pair $(p, q)$, each a $\K$-element (two LE words) \\
\quad \code{u64} + $n_u \times 16$\,B & fold count $n_u$, then \code{us[l]} as LE \code{u128}s \\
\quad \code{Transcribable} & \code{presums[l]} \\
\quad \code{u64} + $n_{sv} \times 16$\,B & ring length $n_{sv}$ ($=128$ when honest), then \code{s\_v} \\
\code{u64} + $n$ B & length-prefixed \textbf{bincode 1.3} blob of \code{lig} (flock's own pinned encoder) \\
\bottomrule
\end{tabular}
\end{center}

\code{Transcribable} values are self-describing: a \code{u32} LE byte
length, then the payload.  A \code{SumcheckProof} payload is, in order:
a 16-byte field-modulus tag (the $\mathtt{0x87}$ pattern); a \code{u32}
message count; per message, a \code{u32} evaluation count followed by
that many 16-byte evaluations; and the 16-byte claimed sum.  A
\code{MultiDegreeSumcheckProof} payload is: the 16-byte modulus tag; a
\code{u32} group count; a \code{u32} round count; one \code{u32} degree
per group; then all groups' round messages (per group, per round, the
degree-many 16-byte evaluations); then one 16-byte claimed sum per group.

\paragraph{What \code{from\_bytes} rejects, exactly.}  The outer reader
returns \code{CodecError::Truncated} whenever a \code{u64}/\code{u128}/%
$\K$-element/blob read would run past the buffer (including overflow of
the cursor), and \code{CodecError::Bincode} when the embedded
\code{LigeritoProof} blob fails bincode decoding.  Everything else is
deferred to verification: counts inconsistent with the statement
($L \ne \lceil q_{\mathrm{bits}}/\cw\rceil$, wrong layer counts, wrong
ring lengths, \dots) decode fine and are rejected by the shape checks of
Section~\ref{sec:schedule}.

\paragraph{Codec caveats (robustness, not soundness).}  Three properties
a reviewer should know; none creates a false-acceptance path, since every
decoded proof still faces the full verifier:
\begin{itemize}[itemsep=1pt]
\item \emph{Panic paths on malformed streams.}  The \code{Transcribable}
  sub-decoders are \code{assert!}/\code{expect}-based: a truncated or
  length-inconsistent inner payload, or an embedded field-modulus tag
  $\ne \mathtt{0x87}$, aborts by panic rather than returning
  \code{CodecError}.  A host feeding untrusted bytes should treat
  \code{from\_bytes} as panicking-on-hostile-input (e.g.\ wrap it) until
  hardened.
\item \emph{Non-strict decoding.}  Trailing bytes after the
  \code{lig} blob are ignored (no end-of-stream check), and any
  \code{sc\_x} presence flag $\ne 1$ decodes as ``absent.''  Hence
  distinct byte streams can decode to one proof; canonicity
  (\emph{re}-serialization being byte-identical, pinned by test) is a
  property of \code{to\_bytes} output, not of arbitrary accepted input.
\item \emph{No semantic validation at decode.}  By design: the codec is a
  framing layer; the verifier is the validator.
\end{itemize}

% ===========================================================================
\section{Verification order and the Fiat--Shamir schedule}
\label{sec:schedule}

\subsection{Transcript primitives}
\label{sec:transcript}

The transcript is a Blake3 hasher (\code{Blake3Transcript},
\code{src/transcript/mod.rs}).  Byte-level semantics:
\begin{itemize}[itemsep=1pt]
\item \code{absorb\_slice(buf)} updates the hasher with
  $\mathtt{0x06} \,\|\, \mathrm{buf} \,\|\, \mathtt{0x07}$.
\item \code{absorb\_gfs(tag, vals)} (\code{src/merged\_forest.rs}) and
  \code{absorb\_gf\_slice(tag, vals)} (\code{src/ligerito.rs}) call
  \code{absorb\_slice} on one domain-tag byte followed by each
  $\K$-element as two LE \code{u64} words.  Tags used
  on this path: $\mathtt{0x30}$ (forest roots), $\mathtt{0x32}$ (layer
  closing pair), $\mathtt{0x20}$ (ring-switch \code{s\_v}).  (Tag
  $\mathtt{0x37}$, virtual-XOR externals, does not occur here.)
\item \code{absorb\_random\_field(v)} — used by the sumcheck layers —
  updates with
  \[
    \mathtt{0x03} \,\|\, \mathrm{modulus} \,\|\, \mathtt{0x05} \,\|\,
    \mathtt{0x01} \,\|\, \mathrm{inner} \,\|\, \mathtt{0x03},
  \]
  where for $\K$ the modulus bytes are the fixed 16-byte LE pattern of
  $\mathtt{0x87}$ (the reduction pentanomial
  $X^{128}+X^7+X^2+X+1$ with the $X^{128}$ bit implicit) and the inner
  bytes are the element's 16-byte LE representation.
\item \code{get\_field\_challenge::<Gf>()} squeezes 16 bytes from the
  Blake3 XOF (\emph{without} advancing the state), then ratchets the state
  with $\mathtt{0x12} \,\|\, \text{squeezed bytes} \,\|\, \mathtt{0x34}$,
  and interprets the squeezed bytes as a $\K$-element (raw 128 bits,
  LSB-first; every 16-byte string is a valid element, so challenges are
  uniform over $\K$).  \code{get\_field\_challenges(n)} is $n$ sequential
  such calls.
\item \code{ZincChallenger} (\code{src/ligerito\_flock.rs}) drives
  flock's \code{Challenger} trait from the same transcript:
  \code{observe\_label}/\code{observe\_bytes}/\code{observe\_f128(\_slice)}
  map to \code{absorb\_slice} (elements as 16-byte LE),
  \code{sample\_f128} maps to \code{get\_field\_challenge}, and the
  proof-of-work hooks squeeze one field element as grinding seed and
  absorb the 8-byte LE nonce via \code{absorb\_slice}
  (\code{verify\_pow} absorbs the nonce whether or not the grinding check
  passes, keeping the transcript in lockstep; an honest verifier rejects
  on failure inside the oracle).  Thus the oracle's sub-schedule
  \emph{continues the same Fiat--Shamir chain}; its contents are out of
  scope.
\end{itemize}

\subsection{The verifier, step by step}
\label{sec:steps}

Below is the complete operational order of
\code{verify\_mle\_eval\_mod\_q\_ligerito}, with every check anchored
\Vc{$\cdot$} and every transcript action marked
$\mathbf{A}$ (absorb) / $\mathbf{S}$ (squeeze).  The prover performs the
identical $\mathbf{A}/\mathbf{S}$ sequence at the corresponding stages, so
transcripts stay in lockstep; any divergence misaligns every later
challenge and the run rejects (with overwhelming probability) at the first
challenge-dependent check.

\begin{enumerate}[leftmargin=2.2em, itemsep=2pt]
\item \emph{Setup (no transcript activity).}  Recompute
  $\cw$, $L$; \Vc{0} reject \code{Shape} unless
  $|\code{mfs}| = |\code{us}| = |\code{presums}| = |\code{rings}| = L$;
  recompute the chunk weights $W^{(l)}$ from \code{row\_weights\_q}; set
  $\mathrm{bound} = 2^{\cw + t + W} = 2^{127}$.

\item \emph{Per chunk $l = 0, \dots, L-1$, sequentially:}
  \begin{enumerate}[leftmargin=1.6em, itemsep=1pt]
  \item \Vc{1} \emph{Range:} for every $c$:
    $u^{(l)}_c < 2^{127}$, else \code{ChunkRange\{chunk: $l$, col: $c$\}}.
    (Runs before any of chunk $l$'s transcript activity.)
  \item Enter \code{verify\_int\_eval\_merged\_common} with weights
    $W^{(l)}$:
    \begin{enumerate}[leftmargin=1.6em, itemsep=1pt]
    \item \Vc{2} $|\code{us[l]}| = 2^s$, else \code{Forest}.
    \item \Vc{3} $\code{is\_generator}(\alpha)$, else
      \code{ChallengeNotGenerator} (Section~\ref{sec:binding}; evaluated
      identically every chunk).
    \item \Vc{4} $\max_c u^{(l)}_c < 2^{128}-1$, else
      \code{Magnitude}.  Vacuous given \Vc{1} (kept for shape parity with
      the unchunked path).
    \item \emph{Recompute roots} $R^{(l)}_c = \alpha^{u^{(l)}_c}$
      (fixed-base comb, \code{FixedBasePow::new($\alpha$, 128, 8)}).
      \textbf{(R)} from proof data — binding comes from the next step's
      absorb.
    \item Run \code{verify\_merged\_forest} (Section~\ref{sec:forest}):
      \Vc{5}–\Vc{11}, absorbs
      $\mathbf{A}(\mathtt{0x30}\!:\!\text{roots})$,
      $\mathbf{S}(\zeta \in \K^s)$, then per layer the phase-A/B sumcheck
      exchanges and
      $\mathbf{A}(\mathtt{0x32}\!:\!\text{pair})$,
      $\mathbf{S}(\lambda_\ell)$.  Output: exit point
      $z^{(l)} = (z_x^{(l)}, z_c^{(l)}) \in \K^{d} \times \K^{s}$ and exit
      value $e_d^{(l)}$.  Failure maps to \code{Forest}.
    \item Run \code{MultiDegreeSumcheck::verify\_as\_subprotocol} on
      \code{presums[l]} with round count $d$
      (Section~\ref{sec:linkage}): absorbs
      $\mathbf{A}(d)$, $\mathbf{A}(n_{\mathrm{groups}})$,
      $\mathbf{A}(\text{degree})$ (each via
      \code{absorb\_random\_field}), then per round
      $\mathbf{A}(\text{round message})$, $\mathbf{S}(\text{challenge})$,
      $\mathbf{A}(\text{challenge})$; checks \Vc{12} (round count $= d$)
      and \Vc{13} (per-round message length $=$ declared degree, enforced
      at the end).  Output: shared point $r^{*(l)} \in \K^{d}$ and
      expected evaluation $E^{(l)}$.  Failure maps to \code{PreSumcheck}.
    \item \Vc{14} $\code{presums[l].claimed\_sums()} = [\,e_d^{(l)} + 1\,]$
      (one group, exactly this value), else \code{PreSumcheck}.
    \item Recompute
      $\widehat{R}^{(l)} = \sum_i q^{(l)}_{\mathrm{rowbit}}[i]\cdot
       \eq(i, r^{*(l)})$ with
      $q^{(l)}_{\mathrm{rowbit}}[i] = \eq(i, z_x^{(l)})\cdot
       (\alpha^{W^{(l)}_b 2^j} - 1)$, $i = (b \ll \log_2 W)|j$
      (\code{row\_bit\_weights}; the verifier's one $O(2^{d})$ step).
      \Vc{15} $\widehat{R}^{(l)} \ne 0$, else \code{RHatZero}.  Set
      $\mu_l = E^{(l)} \cdot (\widehat{R}^{(l)})^{-1}$ and
      $\mathit{pt}_l = (r^{*(l)} \,\|\, z_c^{(l)}) \in \K^{d+s}$.
    \end{enumerate}
  \end{enumerate}

\item \emph{Ring-switch checks, per chunk $l = 0, \dots, L-1$:}
  \Vc{16} $|\code{s\_v}^{(l)}| = 128$, else \code{Shape};
  \Vc{17} $\sum_{v=0}^{127} \eq(v, \mathit{pt}_l[0..7]) \cdot s^{(l)}_v
   = \mu_l$, else \code{RingSwitchClaim}; then
  $\mathbf{A}(\mathtt{0x20}\!:\!\code{s\_v}^{(l)})$.
  (Check first, absorb second; the prover absorbs unconditionally at the
  same position.)

\item $\mathbf{S}(r'' \in \K^{7})$, then
  $\mathbf{S}(\eta \in \K^{L})$ — seven then $L$ sequential field
  challenges.  Build $\eq_{r''} \in \K^{128}$.

\item \emph{Target.}  Recompute
  $\beta_l = \sum_{u=0}^{127} \eq_{r''}[u] \cdot
   \big(\mathsf{transpose}_{128}(\code{s\_v}^{(l)})\big)[u]$
  and $\beta = \sum_l \eta_l \beta_l$
  (\code{transpose\_bits\_128}: the $128\times128$ bit transpose,
  $\mathrm{bit}_v(\mathrm{out}[u]) = \mathrm{bit}_u(\mathrm{in}[v])$).

\item \Vc{18} \emph{The oracle call} (Definition~\ref{def:oracle}):
  \[
    \RecVerify(\mathit{vc},\, \code{lig},\, m_p,\, \beta,\, \mathsf{root},\,
     \mathsf{EvalB},\, \mathcal{C}),
  \]
  with $\mathsf{EvalB}$ the closure over
  $\{\mathit{pt}_l^{\mathrm{hi}}\}_l = \{\mathit{pt}_l[7..]\}_l$,
  $\eq_{r''}$ and $\eta$ described in Section~\ref{sec:linkage}, and
  $\mathcal{C} = \code{ZincChallenger(transcript)}$.  Reject
  \code{LigeritoReject} if the oracle rejects.  All transcript activity
  inside is the oracle's.

\item \Vc{19} \emph{Read-off} (after the oracle; pure recomputation):
  \[
    y_{\mathrm{ver}} \;=\; \sum_{c} w'_c \sum_{l} 2^{\cw l}\, u^{(l)}_c
    \quad\text{in } R
    \qquad (\code{recombine\_read\_off}),
  \]
  reject \code{ReadOff} unless $y_{\mathrm{ver}} = \code{claimed}$.
\end{enumerate}

{\sloppy
Error type: \code{FlockRsError}; the variants reachable here are
\code{RingSwitch} (of \code{Shape} or \code{RingSwitchClaim}),
\code{ChunkRange}, \code{LigeritoReject}, and \code{Common} (of
\code{Forest}, \code{ChallengeNotGenerator}, \code{Magnitude},
\code{PreSumcheck}, \code{RHatZero}, or \code{ReadOff}).\par}
The verifier is
fail-fast: the \emph{first} failed check aborts, so later transcript
activity does not occur on rejection (relevant only to error-behavior
review, not to soundness).

\subsection{The Fiat--Shamir invariant}
\label{sec:fsinvariant}

\emph{Every squeeze follows the absorbs of everything that squeeze must
bind.}  Concretely, in transcript order:
\begin{itemize}[itemsep=1pt]
\item $\zeta^{(l)}$ follows $\mathbf{A}(\mathtt{0x30})$ of chunk $l$'s
  recomputed roots — this is the absorb that binds the sent folds
  \code{us[l]} (the roots are a deterministic function
  $\alpha^{u^{(l)}_c}$ of them), and transitively the entire earlier
  transcript, including all previous chunks.
\item Each sumcheck round challenge follows the absorb of that round's
  message(s); each layer's line challenge $\lambda_\ell$ follows
  $\mathbf{A}(\mathtt{0x32})$ of that layer's closing pair; the
  pre-sumcheck's metadata ($d$, group count, degree) is absorbed before
  its first round.
\item $r''$ and $\eta$ follow the absorbs of \emph{all} $L$ ring
  messages $\code{s\_v}^{(l)}$ (and everything before them); the oracle
  target $\beta$ is a function only of absorbed data
  ($\code{s\_v}$) and of $r'', \eta$ themselves.
\item The oracle's internal challenges continue the same chain, hence
  follow everything above.
\end{itemize}

\begin{obligation}[Statement binding is the caller's]
\label{ob:root}
Neither \code{verify\_mle\_eval\_mod\_q\_ligerito} nor the in-repo tests
absorb the commitment root, the weights, $q_{\mathrm{bits}}$, the claimed
$y$, or $\alpha$ into the transcript: the API takes them as inputs, and
the module documentation (\code{src/ligerito.rs}) places the absorb of the
root ``before any challenge is drawn'' with the \emph{caller}.  In a
Fiat--Shamir deployment the host must (i) absorb the commitment root and
all statement data into this transcript before the first squeeze-bearing
call, and (ii) sample any statement randomness (the evaluation point
behind the weights; $\alpha$, if random) from that transcript.  For
$\alpha$ specifically, binding soundness does not require randomness —
any \emph{generator} works (Section~\ref{sec:binding}) — but it must be
fixed before the prover commits to \code{us}.
\end{obligation}

% ===========================================================================
\section{The merged-forest GKR verifier}
\label{sec:forest}

\code{verify\_merged\_forest(transcript, roots, proof, depth = $d$, s)}
(\code{src/merged\_forest.rs}) verifies that a family of $2^s$ product
trees of depth $d$, whose leaf layer is the (never-materialized)
multilinear family
$\Lambda^{(l)}[c][i] \in \K$, has per-tree roots equal to the supplied
\code{roots}, reducing the statement to a single evaluation claim on the
leaf MLE.  On this path the leaves are the bit-affine values
\begin{equation}\label{eq:leaves}
  \Lambda^{(l)}[c][i] \;=\; 1 + M[c][i]\cdot \tau^{(l)}(i),
  \qquad
  \tau^{(l)}(i) \;=\; \alpha^{W^{(l)}_b 2^j} - 1
  \quad (i = (b \ll \log_2 W)|j),
\end{equation}
so that $\prod_i \Lambda^{(l)}[c][i] = \alpha^{\sum_b W^{(l)}_b
\INT(D(b,c))} = \alpha^{u^{(l)}_c}$ for the \emph{true} fold — the whole
point of the forest.  (In $\K$, $-1 = +1$; we keep the $-1$ form for
readability.)

\paragraph{Roots are derived, not carried.}  The caller recomputes
$R_c = \alpha^{u_c}$ from the sent folds and passes them in; the proof
object has no root field.  The very first transcript action,
$\mathbf{A}(\mathtt{0x30}\!:\!R_0,\dots,R_{2^s-1})$, is what binds the
folds: a prover whose forest messages were generated against different
roots faces challenges derived from \emph{these} roots and fails the layer
checks with overwhelming probability.  \emph{The absorb is the binding.}

\paragraph{Entry.}  \Vc{5} \code{Shape} unless $|\mathrm{roots}| = 2^s$
and $|\code{layers}| = d$.  Then $\mathbf{S}(\zeta \in \K^s)$ and the
verifier computes the entry claim \emph{itself}:
\[
  C_0 \;=\; \code{mle\_at}(\mathrm{roots}, \zeta) \;=\; \mle{R}(\zeta)
  \qquad \textbf{(R)},
\]
by $s$ successive low-bit binds ($O(2^s)$ work).  No prover data is
involved in $C_0$.

\paragraph{Layer loop.}  The trees' levels form one multilinear family
$L_\ell(x, c)$ over $(x, c) \in \bits^\ell \times \bits^s$ ($x$ = in-tree
index, low; $c$ = tree index, high) with the top-bit product recursion
$L_\ell(x,c) = L_{\ell+1}(x, 0, c) \cdot L_{\ell+1}(x, 1, c)$ (the new
child bit is the \emph{top} in-tree bit).  The verifier maintains a
running claim $C_\ell = \mle{L_\ell}(z_x, z_c)$, starting from
$(z_x, z_c) = (\varepsilon, \zeta)$, $C_0$ as above.  For each layer
$\ell = 0, \dots, d-1$, the claim
\[
  C_\ell \;=\; \sum_{x \in \bits^\ell} \sum_{c \in \bits^s}
  \eq(x, z_x)\, \eq(c, z_c)\; L_{\ell+1}(x,0,c)\, L_{\ell+1}(x,1,c)
\]
is discharged in two sumcheck phases plus a line step:
\begin{enumerate}[leftmargin=2.2em, itemsep=1pt]
\item \emph{Phase A} (in-tree variables; only for $\ell \ge 1$):
  \Vc{6} layer $0$ must have \code{sc\_x} absent and layers $\ge 1$
  present (\code{Shape});
  \Vc{7} the phase-A proof's \code{claimed\_sum} must equal the running
  claim $C_\ell$ (\code{LayerClaim}) — this is where the chain from the
  absorbed roots reaches the layer.  Then
  \code{MLSumcheck::verify\_as\_subprotocol(transcript, $\ell$, 3, sc\_x)}
  runs an $\ell$-round, \emph{degree-3} sumcheck (degree pinned by the
  verifier, not the proof), yielding a point $r_x \in \K^\ell$ and an
  expected evaluation $E_A$.
  \Vc{10} $E_A = \eq(r_x, z_x)\cdot \code{sc\_c.claimed\_sum}$
  (\code{LayerClaim}) — the factored-$\eq$ linkage that hands the residual
  sum to phase B.  (At $\ell = 0$, \Vc{7} instead checks
  $\code{sc\_c.claimed\_sum} = C_0$ directly.)
\item \emph{Phase B} (tree-index variables):
  \code{MLSumcheck::verify\_as\_subprotocol(transcript, $s$, 3, sc\_c)}
  yields $r_c \in \K^s$ and $E_B$.
  \Vc{11} $E_B = \eq(r_c, z_c)\cdot p \cdot q$ where
  $(p, q) = \code{pair}$ is the layer's closing child pair
  (\code{LayerClaim}).  The honest values are
  $p = \mle{L_{\ell+1}}(r_x, 0, r_c)$, $q = \mle{L_{\ell+1}}(r_x, 1, r_c)$.
\item \emph{Line step:}
  $\mathbf{A}(\mathtt{0x32}\!:\!p, q)$, $\mathbf{S}(\lambda_\ell)$, and the
  claim advances along the line through the pair:
  \[
    C_{\ell+1} \;=\; p + \lambda_\ell\,(p + q), \qquad
    z_x \leftarrow r_x \,\|\, \lambda_\ell, \qquad
    z_c \leftarrow r_c .
  \]
  ($\lambda_\ell$ is the code's per-layer \code{mu}; renamed here to avoid
  clashing with the residual value $\mu_l$.)
\end{enumerate}
After $d$ layers the exit is $z = (z_x, z_c) \in \K^{d+s}$ with exit value
$e_d = C_d$, honestly $\mle{\Lambda}(z_x, z_c)$ — pinned by the
\code{merged\_forest\_roundtrips} test, which checks the exit value equals
\code{mle\_at(leaves, $z$)}.

\paragraph{Inside the generic sumcheck.}
(\code{src/piop/sumcheck.rs} + \code{verifier.rs}.)  Both phases use the
same subprotocol.  It first absorbs the variable count and the degree
(via \code{absorb\_random\_field}), rejects unless the message count
equals the variable count \Vc{8}, then per round: absorbs the round
message, squeezes the round challenge, and absorbs the challenge back.
Round messages carry only the \emph{tail} evaluations
$P_i(1), \dots, P_i(\deg)$; the verifier reconstructs
$P_i(0) := \mathrm{expected}_{i-1} - P_i(1)$, so the per-round sum
identity $P_i(0) + P_i(1) = \mathrm{expected}_{i-1}$ holds \emph{by
construction} and the protocol's content is the chaining
$\mathrm{expected}_i = P_i(r_i)$ (Lagrange interpolation on
$\{0,\dots,\deg\}$) together with the caller's use of the final expected
evaluation (checks \Vc{10}, \Vc{11} above; check \Vc{14}+division in
Section~\ref{sec:linkage}).  Per-round message length is enforced equal to
the degree at the end \Vc{9} (\code{MaxDegreeExceeded}); the absorbs use
the messages as sent.

\paragraph{Provenance.}  In this section: the roots are \textbf{(R)}
recomputed from proof data and bound by $\mathbf{A}(\mathtt{0x30})$;
$C_0$ is \textbf{(R)} from bound data plus $\zeta$; the round messages,
closing pairs, and \code{claimed\_sum} fields are \textbf{(A)} — absorbed,
and additionally checked (\Vc{7}, \Vc{10}, \Vc{11}) against
already-chained values; all challenges are squeezed; $e_d$ and $z$ are
\textbf{(R)} functions of bound data and challenges.  Nothing here touches
the commitment yet: the forest alone binds \code{us} to \emph{some}
multilinear leaf family; identifying that family with the committed bits
is the next section's job.

% ===========================================================================
\section{The exit-claim linkage: pre-sumcheck, ring switch,
  $\eta$-batching, one oracle call}
\label{sec:linkage}

This is the load-bearing glue: it turns each chunk's forest exit claim —
an evaluation of the \emph{leaf} MLE, which entangles the committed bits
with the public $\alpha$-power table — into a $\K$-linear functional of
the committed bits alone, and hands the batch of $L$ such functionals to
the oracle as \emph{one} inner-product claim.

\subsection{From the exit claim to a bit-MLE claim: the pre-sumcheck}

By \eqref{eq:leaves} and multilinearity of the extension (the MLE of the
all-ones table is $1$):
\begin{equation}\label{eq:exit}
  e_d^{(l)} \;-\; 1
  \;=\; \sum_{i \in \bits^{d}} \sum_{c \in \bits^s}
        \eq(i, z_x^{(l)})\,\eq(c, z_c^{(l)})\; M[c][i]\;\tau^{(l)}(i)
  \;=\; \sum_{i \in \bits^{d}} q^{(l)}_{\mathrm{rowbit}}[i]\cdot m^{(l)}[i],
\end{equation}
where
$q^{(l)}_{\mathrm{rowbit}}[i] = \eq(i, z_x^{(l)})\,\tau^{(l)}(i)$
(\code{row\_bit\_weights} at $\rho = z_x^{(l)}$) and
$m^{(l)}[i] = \sum_c \eq(c, z_c^{(l)})\, M[c][i]$ is the
$z_c$-combined row.  In $\K$, $e_d - 1 = e_d + 1$; the code checks the
$+1$ form.

The pre-sumcheck (\code{presums[l]}) is a $d$-round degree-2 sumcheck on
the right-hand side of \eqref{eq:exit}, verified by
\code{MultiDegreeSumcheck::verify\_as\_subprotocol} with one group.  The
verifier:
\begin{itemize}[itemsep=1pt]
\item checks \Vc{14}
  $\code{claimed\_sums()} = [\,e_d^{(l)} + 1\,]$ — the sumcheck's asserted
  sum is not free data; it must equal the forest-chained exit value, and
  the slice comparison simultaneously pins the group count to one;
\item runs the rounds (shared-challenge machinery, identical
  absorb/squeeze pattern to Section~\ref{sec:forest}'s subprotocol, with
  round count pinned to $d$ by the verifier), obtaining the point
  $r^{*(l)}$ and expected evaluation $E^{(l)}$, honestly
  $\widehat{R}^{(l)} \cdot \mle{m^{(l)}}(r^{*(l)})$;
\item recomputes $\widehat{R}^{(l)} =
  \mle{q^{(l)}_{\mathrm{rowbit}}}(r^{*(l)})$ from public data
  \textbf{(R)} — the $O(2^{d})$ table build plus an inner product with
  the $\eq(\cdot, r^{*(l)})$ table — rejects \Vc{15} if it is zero
  (\code{RHatZero}; for honest runs this happens with probability
  $\approx 2^{-128}$ over the challenges, \emph{except} in the degenerate
  statement where every weight of the chunk is zero, which makes
  $q^{(l)}_{\mathrm{rowbit}} \equiv 0$ and rejects deterministically — a
  completeness corner, not a soundness issue; an adversary gains nothing
  by hoping for a zero since rejection is the outcome), and
  \emph{divides}:
  \[
    \mu_l \;=\; E^{(l)} \cdot \big(\widehat{R}^{(l)}\big)^{-1},
  \]
  which for an honest prover equals
  $\mle{m^{(l)}}(r^{*(l)}) = \mle{M}(r^{*(l)}, z_c^{(l)})$ — a pure
  bit-MLE evaluation at the point
  $\mathit{pt}_l = (r^{*(l)} \| z_c^{(l)})$, with the $\alpha$-power
  weight factor stripped.  The division is the ``de-black-boxing'':
  everything $\alpha$-dependent is now on the verifier's side of the
  equation.
\end{itemize}

\begin{remark}[The pre-sumcheck degree rides the proof]
\label{rem:degree}
Unlike the forest layers (degree pinned to $3$ by the verifier), the
pre-sumcheck's per-group degree is read from the proof, absorbed, and then
enforced per round.  A malicious prover may declare any degree
$\delta < 2^{32}$ (the codec's \code{u32}); the sumcheck's per-round
soundness error grows linearly, $\delta / |\K|$, i.e.\ at most
$2^{32-128}$ per round — negligible.  The declared degree is absorbed
before any round challenge, so it cannot be chosen adaptively per round.
This is a documented asymmetry, not a gap.
\end{remark}

\subsection{From $L$ bit-MLE claims to one oracle call: the ring switch}

Each $\mathit{pt}_l \in \K^{d+s}$ splits at the packing boundary:
$r_{\mathrm{lo}}^{(l)} = \mathit{pt}_l[0..7]$ (the seven in-pack row-bit
coordinates) and $\mathit{pt}_l^{\mathrm{hi}} = \mathit{pt}_l[7..] \in
\K^{m_p}$.  Writing $\mle{M}_v$ for the MLE of the $v$-th bit-slice of the
packed vector $P$ (i.e.\ $P[y] = \sum_v \mle{M}_v$-table$[y]\cdot\beta_v$),
multilinearity gives
$\mle{M}(\mathit{pt}_l) = \sum_{v=0}^{127}
 \eq(v, r_{\mathrm{lo}}^{(l)})\cdot \mle{M}_v(\mathit{pt}_l^{\mathrm{hi}})$.
The prover sends, per chunk, the 128 partial evaluations
$s^{(l)}_v$ claimed to be $\mle{M}_v(\mathit{pt}_l^{\mathrm{hi}})$
(\code{RingSwitchProof.s\_v}); the verifier:
\begin{itemize}[itemsep=1pt]
\item \Vc{17} checks the one $\K$-linear consistency equation
  $\sum_v \eq(v, r_{\mathrm{lo}}^{(l)})\, s^{(l)}_v = \mu_l$ — this ties
  the message to the forest-chained value $\mu_l$ (one constraint out of
  128 degrees of freedom; the rest are bound below);
\item absorbs each message
  ($\mathbf{A}(\mathtt{0x20})$, all $L$ of them in chunk order), then
  squeezes $r'' \in \K^7$ and $\eta \in \K^L$;
\item recombines each message through the $\Ftwo$-linear batching map
  $\Phi_{r''}\colon \beta_u \mapsto \eq_{r''}[u]$: with
  $s_u^{(l)} = \mathsf{transpose}_{128}(s^{(l)})$ (an exact basis
  reindexing, \code{transpose\_bits\_128}),
  \[
    \beta_l \;=\; \sum_{u=0}^{127} \eq_{r''}[u]\; s_u^{(l)},
    \qquad
    \beta \;=\; \sum_{l} \eta_l\, \beta_l \qquad \textbf{(R)} .
  \]
  For honest messages, $\beta_l = \sum_y P[y]\cdot B_l(y)$ with
  \[
    B_l(y) \;=\; \Phi_{r''}\!\big(\eq(\mathit{pt}_l^{\mathrm{hi}},
    y)\text{-table}[y]\big),
  \]
  the ring-switch recombination identity (Flock, App.~B; ring-switching
  due to Diamond--Posen).  If some $s^{(l)}$ differs from the true partial
  evaluations, then over the choice of $r''$ (squeezed \emph{after} the
  absorb of $s^{(l)}$) the recombined $\beta_l$ differs from
  $\sum_y P[y] B_l(y)$ except with probability at most $7 \cdot 2^{-128}$
  (Schwartz--Zippel; the difference is a nonzero multilinear polynomial in
  the $7$ coordinates of $r''$), and over $\eta$ the batch value $\beta$
  then differs from $\sum_y P[y] B(y)$ except with probability
  $2^{-128}$ (the difference is a nonzero linear form in $\eta$) —
  leaving a false $(\beta, B)$ pair to the oracle, which rejects by
  Definition~\ref{def:oracle}.
\end{itemize}

\paragraph{The oracle call.}  The combined weight function is
$B(y) = \sum_l \eta_l\, B_l(y)$.  The verifier never materializes $B$
(that would cost $O(2^{m_p})$); instead it hands the oracle the
$\mathsf{EvalB}$ callback, which evaluates the MLE $\mle{B}$ at
oracle-chosen points succinctly:
\[
  \mathsf{EvalB}(\rho, k')[y] \;=\; \sum_l \eta_l \cdot
  \code{residual\_b\_evals}(\rho,\, k',\, \mathit{pt}_l^{\mathrm{hi}},\,
  \eq_{r''})[y],
\]
the tensor-algebra evaluation of
$\mle{\Phi_{r''} \circ \eq(\mathit{pt}_l^{\mathrm{hi}},\cdot)}$ at
$\rho \,\|\, \mathsf{bits}(y)$ in $O((|\rho| + 2^{k'})\cdot 128^2)$
$\K$-operations, with no $2^{m_p}$-sized table
(\code{residual\_b\_evals} / \code{tensor\_eq\_phi\_eval},
\code{src/ligerito.rs}; pinned against brute force by tests).  The
callback is verifier-side code computing a public function of
$(\mathit{pt}_l^{\mathrm{hi}}, r'', \eta)$ and the oracle's own
challenges; \emph{when} the oracle invokes it is inside the box.  The
oracle call \Vc{18} then certifies $\sum_y P[y] B(y) = \beta$ against the
committed root — the only step in the entire verification that touches
the commitment.

\subsection{Provenance of every verifier-side quantity}

\begin{table}[ht]
\centering\footnotesize
\begin{tabular}{p{0.38\textwidth}cp{0.44\textwidth}}
\toprule
Quantity & Class & Bound / checked by \\
\midrule
$\cw$, $L$, chunk weights $W^{(l)}$, bound $2^{127}$ & (R) & statement ($p$, $q_{\mathrm{bits}}$, $w$) \\
\code{us[l]} (sent folds) & (A) & \Vc{1} range; absorbed as roots $\alpha^{u}$ (tag $\mathtt{0x30}$); \Vc{19} \\
roots $R^{(l)}_c$ & (R) & deterministic in \code{us[l]}, $\alpha$; the absorb binds \\
$\zeta^{(l)}$, $\lambda_\ell$, round challenges, $r^{*(l)}$, $r''$, $\eta$ & — & squeezed (uniform over $\K$) \\
entry claim $C_0^{(l)} = \mle{R}(\zeta)$ & (R) & \code{mle\_at} on bound roots \\
sumcheck round messages (all layers, presum) & (A) & absorbed; length/degree checks \Vc{8}\Vc{9}\Vc{12}\Vc{13} \\
\code{claimed\_sum} fields (\code{sc\_x}, \code{sc\_c}) & (A) & \Vc{7} = running claim; \Vc{10} phase linkage \\
closing pairs $(p,q)$ & (A) & absorbed (tag $\mathtt{0x32}$); \Vc{11} against $E_B$ \\
exit $(z^{(l)}, e_d^{(l)})$ & (R) & chained from bound data + challenges \\
presum \code{claimed\_sums} & (A) & \Vc{14} $= [e_d^{(l)}+1]$ \\
presum declared degree & (A) & absorbed; per-round length \Vc{13} (Remark~\ref{rem:degree}) \\
$\widehat{R}^{(l)}$, $\mu_l$, $\mathit{pt}_l$ & (R) & public table + division; \Vc{15} \\
$\code{s\_v}^{(l)}$ & (A) & \Vc{17} one linear check vs.\ $\mu_l$; absorbed (tag $\mathtt{0x20}$) \\
$\beta_l$, $\beta$ & (R) & deterministic in bound $\code{s\_v}$, $r''$, $\eta$ \\
$\mathsf{EvalB}$ values & (R) & public function of points, $r''$, $\eta$, oracle challenges \\
$\sum_y P[y] B(y) = \beta$ & (O) & \Vc{18}, Definition~\ref{def:oracle} \\
$y_{\mathrm{ver}}$ & (R) & from \code{us}, $w'$; \Vc{19} vs.\ \code{claimed} \\
$\mathit{vc}$, $m_p$ & (R) & shape-derived (Section~\ref{sec:config}) \\
\bottomrule
\end{tabular}
\caption{Provenance of verifier-side quantities.  (R) recomputed from
public data and verifier randomness; (A) prover data bound by an absorb
and/or an equality check against bound data; (O) oracle-certified.}
\label{tab:provenance}
\end{table}

Reading the chain backwards from the oracle: \Vc{18} certifies the
committed bits enter $\beta$; $\beta$ deterministically encodes the
absorbed $\code{s\_v}$ under post-absorb randomness $(r'', \eta)$;
\Vc{17} ties $\code{s\_v}$ to $\mu_l$; $\mu_l$ is a public-data quotient
of the pre-sumcheck's expected evaluation, whose asserted sum \Vc{14} is
the forest exit value; the forest chains that value to the absorbed
roots under fresh challenges; the roots are $\alpha^{\code{us}}$; and
\Vc{1} + the generator property make the exponent map injective there, so
\code{us} must be the true folds of the committed bits.  \Vc{19} then
evaluates the statement on \code{us}.  Each arrow is one of the checks
above; none is ``the verifier recomputed a value the prover chose
freely.''

% ===========================================================================
\section{Integer binding in the clear}
\label{sec:binding}

This section is the $\K$-side trust anchor: it is what lets the verifier
treat the \emph{sent integers} \code{us} as the true folds once the forest
and the oracle have done their work.

\paragraph{The generator check.}  \Vc{3}
$\code{is\_generator}(\alpha)$ (\code{src/pcs.rs}): reject unless
$\alpha \notin \{0, 1\}$ and
$\alpha^{(2^{128}-1)/p} \ne 1$ for every prime
$p \mid 2^{128}-1$, using the hard-coded squarefree factorization
\[
  2^{128}-1 \;=\; 3 \cdot 5 \cdot 17 \cdot 257 \cdot 641 \cdot 65537
  \cdot 274177 \cdot 6700417 \cdot 67280421310721
\]
(\code{GF128\_ORDER\_PRIME\_FACTORS}: the five Fermat primes dividing
$2^{32}-1$, then $2^{32}+1 = 641 \cdot 6700417$ and
$2^{64}+1 = 274177 \cdot 67280421310721$).  This is the standard
primitive-element test; it passes iff
$\mathrm{ord}(\alpha) = 2^{128}-1 = $ \code{GF128\_MULT\_ORDER}
($= \code{u128::MAX}$).  \emph{Why it is needed:} the exponent map
$n \mapsto \alpha^n$ is injective on $[0, \mathrm{ord}\,\alpha)$; since
$2^{128}-1$ is composite, a non-generator $\alpha$ could have order as
small as $3$, collapsing the binding entirely.  With a generator, the map
is injective on all of $[0, 2^{128}-1)$, which strictly contains the
range-checked domain $[0, 2^{127})$.  About $\prod_p (1 - 1/p) \approx
49\%$ of $\Kx$ are generators, so a randomly sampled $\alpha$ passes after
$\approx 2$ resamples; the verifier simply rejects a non-generator.  The
check is repeated (identically) per chunk; it depends only on $\alpha$.

\paragraph{The range check.}  \Vc{1}, in
\code{verify\_mle\_eval\_mod\_q\_ligerito} itself: for every chunk $l$ and
column $c$, reject \code{ChunkRange\{chunk, col\}} unless
\[
  u^{(l)}_c \;<\; 2^{\,\cw + t + W} \;=\; 2^{127}
\]
(the exponent $\cw + t + W = 127$ by construction of $\cw$; the code
computes the bound generically as
\code{1u128 << (c\_w + t + word\_bits)}).  The bound is exactly the
honest-fold envelope \eqref{eq:chunkfold}: honest chunk folds always
pass, and any sent fold that passes lies in the injective range.  The
check runs \emph{before} the chunk's transcript activity and covers every
entry of \code{us[l]}.  (The inner \Vc{4} magnitude guard,
$\max u < 2^{128}-1$, is inherited from the unchunked path and is
strictly weaker; it is vacuous given \Vc{1}.)

\paragraph{The binding argument.}  Fix a chunk $l$ and column $c$.  Let
$u^{\mathrm{true}} = \sum_b W^{(l)}_b\,\INT(D(b,c))$ be the fold of the
\emph{committed} bits.  The forest + pre-sumcheck + ring switch + oracle
chain (Sections~\ref{sec:forest}--\ref{sec:linkage}) forces — up to the
stated soundness errors — the per-tree product of the leaf family
\eqref{eq:leaves} built from the committed bits to equal the absorbed
root, i.e.
$\alpha^{u^{\mathrm{true}}} = \alpha^{u^{(l)}_c}$.
Both exponents lie in $[0, 2^{127})$: the left by
\eqref{eq:chunkfold}, the right by \Vc{1}.  Injectivity of
$n \mapsto \alpha^n$ on $[0, 2^{128}-1) \supset [0, 2^{127})$ (the
generator property, \Vc{3}) gives $u^{(l)}_c = u^{\mathrm{true}}$.  Both
checks are load-bearing: without \Vc{3} the map need not be injective;
without \Vc{1} the prover could send $u + \mathrm{ord}(\alpha) \cdot k$.
Note the argument needs no randomness in $\alpha$ — any generator fixed
before the prover commits to \code{us} works (cf.\
Obligation~\ref{ob:root}).

\paragraph{Root recomputation.}  $R^{(l)}_c = \alpha^{u^{(l)}_c}$ is
computed with an 8-bit-window fixed-base comb —
\code{FixedBasePow::new} at $(\alpha, 128, 8)$ — a pure optimization of
\code{gf\_pow}: $16$ table lookups/multiplies per exponent instead of a
$128$-step square-and-multiply chain.

\paragraph{The mod-$q$ recombination.}  \Vc{19}
(\code{recombine\_read\_off}, \code{src/pcs.rs}): in the evaluation ring
$R$,
\[
  y_{\mathrm{ver}} \;=\; \sum_{c \in \bits^s} w'_c \cdot
  \sum_{l=0}^{L-1} \big(2^{\cw}\big)^{l} \cdot u^{(l)}_c ,
\]
computed with $2^{\cw}$ lifted once ($\cw \le 126$, so it is an exact
\code{u128}) and powered in $R$; the flattened index is
$\code{v\_flat}[\,l \cdot 2^s + c\,] = \code{us[l][c]}$.  The wide value
$\sum_l 2^{\cw l} u^{(l)}_c$ never materializes as an integer — the whole
point of chunking — and the column weights enter at full width,
unchunked.  Reject \code{ReadOff} unless $y_{\mathrm{ver}} =
\code{claimed}$.  Correctness of the recombination against
\eqref{eq:claim} is the telescoping
$\sum_l 2^{\cw l} W^{(l)}_b = w_b$ (exact under the width precondition)
and the ring homomorphism $\code{R::from} \colon [0, 2^{127}) \to R$
being faithful on the values that occur — a statement-side precondition
on $R$ (for $R = \Fq$ with $q$ a 100-bit prime, reduction mod $q$ is the
canonical map; ``faithful'' means the canonical composite, not
injectivity).

% ===========================================================================
\section{Config derivation and the audited regime}
\label{sec:config}

The Ligerito verifier config $\mathit{vc}$ is \emph{derived from the
instance shape, never from the proof}:
\code{sha\_lig\_configs($m_p$)} (\code{src/ligerito\_flock.rs}) selects,
with $m = m_p + 7$:
\begin{itemize}[itemsep=1pt]
\item $m \ge 22$: the \emph{embedded \textsc{fast} profile} at $m$ — the
  \emph{audited} regime;
\item $m < 22$: \code{ligerito::default\_config} at
  $(m_p,\ \code{log\_batch}=2,\ \code{log\_inv\_rate}=2)$ — ad-hoc,
  \emph{unaudited}, test-only.
\end{itemize}
The embedded profiles are flock's own audited security configs
(deserialized from embedded TOML via \code{LigeritoSecurityConfig},
existing for $m = 22..35$); they also dictate the \emph{commit} shape
(\code{log\_inv\_rate} $=$ \code{log\_inv\_rates[0]},
\code{log\_batch} $=$ \code{initial\_k}), so prover, committer, and
verifier all derive the same parameters from $m_p$ alone.  The ad-hoc
branch builds unique-decoding-regime query counts with no grinding and no
out-of-domain samples, and its per-level parameters are not audited.

\textbf{Scope of this note's claims.}  Everything this note asserts about
the oracle (Definition~\ref{def:oracle}) is asserted \emph{for the
embedded-profile regime $m \ge 22$ only} — every deployed shape in the
README's reference table from $n = 22$ up.  Runs below that boundary
(including the repository's own unit tests, Section~\ref{sec:tests}) use
the unaudited ad-hoc config: the verification \emph{logic} is identical,
but the oracle's concrete soundness there carries no audit and is outside
this note's claims.  A caller may also bypass \code{sha\_lig\_configs}
and pass any \code{LigConfig} to \code{lig\_configs}; the same scoping
applies.  (Performance is not this note's subject, but the boundary is
visible in the README's tables as the $n = 20 \to 22$ step in proof size;
hard-coding the ad-hoc config at large shapes is also catastrophically
slow — measured 292\,s vs.\ tens of ms for an $n=28$ commit.)

% ===========================================================================
\section{What the tests pin}
\label{sec:tests}

The suite (\code{cargo test --release}: 68 lib tests, of which one is an
ignored measurement baseline, plus 2 doctests) pins the verifier's
behavior at these points; test names are exact.

\begin{itemize}[itemsep=2pt]
\item \code{mle\_eval\_mod\_q\_ligerito\_roundtrips}
  (\code{src/ligerito\_flock.rs}) — the end-to-end reference, at
  $(t,s,W) = (10,5,1)$ ($L=1$) and $(4,8,32)$ ($L=2$), with a concrete
  $\Fq$, $q = 2^{100}-15$.  Arms: honest proof accepts; wrong claimed $y$
  rejects with \code{ReadOff} \Vc{19}; a tampered chunk fold
  ($\code{us[0][0]} \leftarrow 2^{128}-2$) rejects with \code{ChunkRange}
  \Vc{1}; a non-generator $\alpha = 1$ rejects with
  \code{ChallengeNotGenerator} \Vc{3}.
\item \code{mod\_q\_ligerito\_proof\_serialization\_roundtrips} — decode
  $\circ$ encode preserves acceptance; re-serialization is byte-identical
  (canonicity of \code{to\_bytes}); flipping a single byte at each of four
  positions ($0$, $\lfloor n/3 \rfloor$, $\lfloor n/2 \rfloor$, $n-1$) is
  rejected — the stream fails to decode or the decoded proof fails
  verification.
\item \code{merged\_forest\_roundtrips} (\code{src/merged\_forest.rs}) —
  prover/verifier exit-claim agreement over four $(d, s)$ shapes; roots
  equal the per-tree leaf products; \emph{the exit value equals
  \code{mle\_at(leaves, $z$)}} (the semantic anchor for
  Section~\ref{sec:forest}); tamper arms: a phase-A \code{claimed\_sum}
  \Vc{7}, a layer-0 closing pair \Vc{11}, and the last layer's phase-B
  \code{claimed\_sum} each reject.
\item \code{lazy\_matches\_eager} and \code{lazy\_multi\_matches\_eager}
  — both lazy prover schedules (L/4 default and
  \code{F2\_FOREST\_SCHEDULE=l8}) are byte-identical to the eager
  reference prover: same roots, same per-layer messages and pairs, same
  exit claim, \emph{same post-run transcript state} (a squeezed challenge
  compares equal).  This is what lets the verifier be schedule-agnostic.
\item \code{tensor\_eval\_matches\_bruteforce\_and\_table} and
  \code{residual\_b\_evals\_matches\_pointwise}
  (\code{src/ligerito.rs}) — the succinct $\mathsf{EvalB}$ algebra equals
  the brute-force $\sum_y \eq(y,\cdot)\,\Phi_{r''}(\cdot)$ and the
  per-boolean-tail pointwise evaluations.
\item \code{field\_bridge\_is\_bit\_identical}
  (\code{src/ligerito\_flock.rs}) — the zinc/flock $\K$ representations
  agree as raw words under $+$ and $\times$, and \code{LOG\_PACKING}
  equals flock's \code{pcs::pack::LOG\_PACKING}.
  \code{neon\_mul\_matches\_scalar\_pipeline} pins the aarch64 NEON field
  pipeline bit-identical to the portable scalar one.
\item \code{packed\_cols\_row\_roundtrip} and the \code{poly/} field-axiom
  and lift/projection tests — layout inversions and $\K$ arithmetic.
\end{itemize}

\paragraph{What the tests do \emph{not} cover.}  Honest gaps a reviewer
should know:
\begin{itemize}[itemsep=2pt]
\item \emph{Fiat--Shamir order mutations.}  No test mutates an absorb
  order, a domain tag, or a framing byte and asserts rejection; the
  schedule of Section~\ref{sec:schedule} is pinned only indirectly
  (honest-roundtrip acceptance requires prover/verifier lockstep, and the
  byte-identity tests pin the two prover variants to each other).  A
  transcript-schedule regression that changed prover and verifier
  \emph{together} would pass the suite.
\item \emph{Codec field coverage.}  The tamper test flips four fixed byte
  positions of one proof; there is no exhaustive per-field tamper sweep,
  and no test distinguishes clean \code{CodecError} returns from the
  panic paths of Section~\ref{sec:serialization}.
\item \emph{Rejection-arm coverage.}  \code{Shape},
  \code{RingSwitchClaim} \Vc{17}, \code{PreSumcheck} \Vc{14},
  \code{RHatZero} \Vc{15}, and \code{LigeritoReject} \Vc{18} have no
  dedicated negative tests in this repository (some are hit
  probabilistically by the byte-flip test).
\item \emph{The audited regime is not unit-tested.}  All
  \code{cargo test} shapes sit below the $m \ge 22$ boundary and run the
  ad-hoc config; the embedded profiles are exercised only by the
  \code{reference\_measure} example and the bench suite
  (Section~\ref{sec:config}).
\item \emph{Batched and virtual-XOR variants} have no tests in this
  repository (they were dropped in the extraction; the code is a verbatim
  port).
\end{itemize}

% ===========================================================================
\section{Trust boundary and caveats}
\label{sec:trust}

\paragraph{The oracle assumption.}  Everything rests on
Definition~\ref{def:oracle} holding for
\code{flock-core}'s recursive Ligerito verifier under the shape-derived
config (Section~\ref{sec:config}).  \code{flock-core} is consumed as a
\emph{local path dependency}
(\code{flock-core = \{ path = ".../flock/crates/flock-core" \}} in
\code{Cargo.toml}, mirroring the upstream pin; it pins \code{bincode 1.3}
and \code{serde 1}; MIT OR Apache-2.0).  There is no version or content
pin beyond the checkout itself — reproducing a verification requires the
same flock tree.  The oracle's internal Fiat--Shamir sub-schedule runs
over the shared transcript through \code{ZincChallenger}
(Section~\ref{sec:transcript}) and is trusted along with the rest of the
box.

\paragraph{Known upstream item: Merkle domain separation.}  flock's
\code{merkle} module does not domain-separate leaf hashing from
internal-node hashing, and its own module note flags it as a
``micro-benchmark module, not production code.''  F2Z inherits flock's
commitment and Merkle verbatim; until flock ships the fix, the
commitment binding — and therefore Definition~\ref{def:oracle} — should
be treated as pre-production.  (Recorded in \code{README.md} and
\code{docs/DESIGN.md}; this note adds nothing beyond scoping its claims
accordingly.)

\paragraph{Caller obligations (summary).}  Obligation~\ref{ob:root}
(statement/root absorbed before challenges; $\alpha$ fixed
pre-commitment); weight vectors of the right lengths ($2^t$, $2^s$) with
$w_b < 2^{q_{\mathrm{bits}}}$; $q_{\mathrm{bits}}$ consistent with the
intended $\Fq$; $R$ of characteristic $\ne 2$ with \code{From<u128>} the
canonical map (a characteristic-2 $R$ would collapse the read-off to
parity — the very obstruction the exponent fold exists to avoid); the
commitment produced at the same shape \code{p} handed to the verifier.
The proof-stream decoder is not hardened against hostile bytes
(Section~\ref{sec:serialization}).

\paragraph{Not zero-knowledge.}  The Ligerito opening reveals queried
committed data; no masking exists anywhere on this path.

\paragraph{Out-of-scope public API.}  The batched multi-polynomial
opener (\code{verify\_rs\_ligerito\_batch}) and the virtual-XOR variants
(\code{verify\_mle\_eval\_mod\_q\_ligerito\_with\_virtual\_xors} /
\code{\ldots\_claims\_only}) share this note's machinery but add
statement-level structure this note does not review.  One contractual
point deserves flagging even here: the virtual-XOR verifiers
\emph{return} a \code{Vec<VirtualXorObligation>} — external-term
residual claims $\hat{e}(\mathit{pt}) = \mathrm{value}$ that ride the
proof \emph{unverified} and that the \emph{caller} must discharge by
other means.  A host that ignores the returned obligations has not
verified those claims.  These variants also carry no tests in this
repository (Section~\ref{sec:tests}).

\paragraph{What a correctness argument for F2Z still owes.}  Beyond this
note's line-by-line checkable description: (i) the Ligerito oracle
itself (proximity soundness, Merkle binding incl.\ the domain-separation
item, its FS sub-schedule) in the embedded-profile regime — and
separately anything run under the ad-hoc config; (ii) the host-level
statement binding of Obligation~\ref{ob:root}; (iii) quantitative
composition of the stated per-step soundness errors (forest layers,
pre-sumcheck, ring-switch recombination, $\eta$-batching) into one bound
— this note states each check's role and error shape but proves nothing;
(iv) hardening of \code{from\_bytes} for hostile inputs.

\paragraph{Verifier cost.}  For the record (README reference table,
Apple M4): verification is milliseconds-class across the deployed range —
$1.2$--$14.8$\,ms for $n = 16$--$32$ at $t \approx 0.6n$, $W=1$, one
mod-$q$ claim — with proofs 43--574\,KiB.  The verifier's only
super-logarithmic work is the per-chunk $O(2^{d})$
$q_{\mathrm{rowbit}}$ table (dominated by $2^t$ fixed-base
exponentiations) and the $O(2^s)$ root/entry/read-off loops.

\begin{thebibliography}{9}\small

\bibitem{soukhanov}
L.~Soukhanov.
\emph{Characteristic-2 field switching} (the char2-fieldswitch note),
\S9: a one-round, characteristic-2 Brakedown instantiation.
The construction F2Z implements.

\bibitem{flock}
B.~B\"unz, J.~Rothblum, R.~Wang (Succinct Labs).
\emph{Flock: Fast Proving for Batch Boolean Computations.}
Source of the ring-switch implementation shape (App.~B) and of
\code{flock-core}, this crate's opener backend.

\bibitem{dp24}
B.~E.~Diamond, J.~Posen.
\emph{Polylogarithmic Proofs for Multilinears over Binary Towers}
(ring-switching).  The packing/recombination technique behind
Section~\ref{sec:linkage}.

\bibitem{ligerito}
A.~Novakovic, G.~Angeris.
\emph{Ligerito: A Small and Concretely Fast Polynomial Commitment
Scheme.}  The recursive opener behind Definition~\ref{def:oracle}.

\bibitem{upstream}
The \code{zinc-plus} repository,
\code{documentation/f2-int-eval-doc/} (branch \code{f2-int-sha}):
the as-deployed long-form exposition this note's notation follows.
Where the deployment and this repository disagree, this repository's
code is authoritative here.

\end{thebibliography}

\end{document}
